\ifdefined\gkver\else\def\gkver{student}\fi
\documentclass[\gkver]{gaokaozhenti}
\begin{document}

\chapter{2023年全国甲卷}
%% paperType: 理综物理部分

\begin{enumerate}
\item
%% number: 1
%% paperName: 2023年全国高考甲卷第 1 题 6 分
%% typeId: 1
%% type: 单选题
%% chapter: 曲线运动
%% point: 平抛运动
%% method: 无
%% score: 6
%% degree: 850
%% duplicateId: 0
%% body:
一同学将铅球水平推出，不计空气阻力和转动的影响，铅球在平抛运动过程中\xzanswer{B}
\fourchoices[answer=B]
{机械能一直增加}
{加速度保持不变}
{速度大小保持不变}
{被推出后瞬间动能最大}
%% answer: B
%% memo:
\memoanswer{
A．铅球做平抛运动，仅受重力，故机械能守恒，A 错误；

B．铅球的加速度恒为重力加速度保持不变，B 正确；

CD．铅球做平抛运动，水平方向速度不变，竖直方向做匀加速直线运动，根据运动的合成可知铅球速度变大，则动能越来越大，CD 错误。

故选 B。
}

\item
%% number: 2
%% paperName: 2023年全国高考甲卷第 2 题 6 分
%% typeId: 1
%% type: 单选题
%% chapter: 原子和原子核
%% point: 人工核反应
%% method: 无
%% score: 6
%% degree: 750
%% duplicateId: 0
%% body:
在下列两个核反应方程中 X + $^{14}_{7}$N $\to$ Y + $^{17}_{8}$O、Y + $^{7}_{3}$Li $\to$ 2X，X 和 Y 代表两种不同的原子核，以 $Z$ 和 $A$ 分别表示 X 的电荷数和质量数，则\xzanswer{D}
\fourchoices[answer=D]
{$Z$ = 1，$A$ = 1}
{$Z$ = 1，$A$ = 2}
{$Z$ = 2，$A$ = 3}
{$Z$ = 2，$A$ = 4}
%% answer: D
%% memo:
\memoanswer{
设 Y 电荷数和质量数分别为 $m$ 和 $n$，根据核反应方程质量数和电荷数守恒可知第一个核反应方程的核电荷数和质量数满足

$A$ + 14 = $n$ + 17，$Z$ + 7 = m + 8

第二个核反应方程的核电荷数和质量数满足

$n$ + 7 = 2$A$，$m$ + 3 = 2$Z$

联立解得 $Z$ = 2，$A$ = 4

故选 D。
}

\item
%% number: 3
%% paperName: 2023年全国高考甲卷第 3 题 6 分
%% typeId: 1
%% type: 单选题
%% chapter: 直线运动
%% point: 运动图象
%% method: 无
%% score: 6
%% degree: 650
%% duplicateId: 0
%% body:
一小车沿直线运动，从 $t$ =  0 开始由静止匀加速至 $t = t_{1}$ 时刻，此后做匀减速运动，到 $t = t_{2}$ 时刻速度降为零。在下列小车位移 $x$ 与时间 $t$ 的关系曲线中，可能正确的是\xzanswer{D}
\onepicture{\includesvg{figs/03}}
%% answer: D
%% memo:
\memoanswer{
$x - t $图像的斜率表示速度，小车先做匀加速运动，因此速度变大即 0 \textasciitilde{} $t_{1}$ 图像斜率变大，$t_{1}$ \textasciitilde{} $t_{2}$ 做匀减速运动则图像的斜率变小，在 $t_{2}$ 时刻停止图像的斜率变为零。

故选 D。
}

\item
%% number: 4
%% paperName: 2023年全国高考甲卷第 4 题 6 分
%% typeId: 1
%% type: 单选题
%% chapter: 曲线运动
%% point: 向心加速度向心力
%% method: 无
%% score: 6
%% degree: 650
%% duplicateId: 0
%% body:
一质点做匀速圆周运动，若其所受合力的大小与轨道半径的 $n$ 次方成正比，运动周期与轨道半径成反比，则 $n$ 等于\xzanswer{C}
\fourchoices[answer=C]
{1}
{2}
{3}
{4}
%% answer: C
%% memo:
\memoanswer{
质点做匀速圆周运动，根据题意设周期

 $T=\frac{k}{r}$

合外力等于向心力，根据

 $F_{\text{合}}=F_{n}=m\frac{4\pi^{2}}{T^{2}}r$

联立可得

 $F_{n}=\frac{4m\pi^{2}}{k^{2}}r^{3}$

其中  $\frac{4m\pi^{2}}{k^{2}}$ 为常数，r 的指数为 3，故题中 n=3。

故选 C。
}

\item
%% number: 5
%% paperName: 2023年全国高考甲卷第 5 题 6 分
%% typeId: 1
%% type: 单选题
%% chapter: 电场
%% point: 带电粒子的轨迹类问题
%% method: 无
%% score: 6
%% degree: 650
%% duplicateId: 0
%% body:
在一些电子显示设备中，让阴极发射的电子束通过适当的非匀强电场，可以使发散的电子束聚集。下列 4 幅图中带箭头的实线表示电场线，如果用虚线表示电子可能的运动轨迹，其中正确的是\xzanswer{A}
\onepicture{\includesvg{figs/05}}
%% answer: A
%% memo:
\memoanswer{
A．电子做曲线运动满足合力指向轨迹凹侧，A 正确；

B．电子做曲线运动满足合力指向轨迹凹侧，对电子受力分析有

可见与电场力的受力特点相互矛盾，B 错误；

C. 电子做曲线运动满足合力指向轨迹凹侧，对电子受力分析有

可见与电场力的受力特点相互矛盾，C 错误；

D．电子做曲线运动满足合力指向轨迹凹侧，对电子受力分析有

可见与电场力的受力特点相互矛盾，D 错误；

故选 A。
}

\item
%% number: 6
%% paperName: 2023年全国高考甲卷第 6 题 6 分
%% typeId: 2
%% type: 多选题
%% chapter: 运动和力
%% point: 牛顿定律中的图象
%% method: 无
%% score: 6
%% degree: 650
%% duplicateId: 0
%% body:
用水平拉力使质量分别为 $m_{\text{甲}}$、$m_{\text{乙} }$的甲、乙两物体在水平桌面上由静止开始沿直线运动，两物体与桌面间的动摩擦因数分别为 $\mu_{\text{甲}}$ 和 $\mu_{\text{乙}}$。甲、乙两物体运动后，所受拉力 $F$ 与其加速度 $a$ 的关系图线如图所示。由图可知\xzanswer{BC}
\onepicture{figs/06.png}
\fourchoices[answer=BC]
{$m_{\text{甲}} < m_{\text{乙}}$}
{$m_{\text{甲}} > m_{\text{乙}}$}
{$\mu_{\text{甲}} < \mu_{\text{乙}}$}
{$\mu_{\text{甲}} > \mu_{\text{乙}}$}
%% answer: BC
%% memo:
\memoanswer{
根据牛顿第二定律有    $F-\mu mg = ma$

整理后有        $F = ma+\mu mg$

则可知 $F - a $图像的斜率为 $m$，纵截距为 $\mu mg$，则由题图可看出

$m_{\text{甲}} > m_{\text{乙}}$，$\mu_{\text{甲}}m_{\text{甲}}g = \mu_{\text{乙}}m_{\text{乙}}g$

则    $\mu_{\text{甲}} < \mu_{\text{乙}}$

故选 BC。
}

\item
%% number: 7
%% paperName: 2023年全国高考甲卷第 7 题 6 分
%% typeId: 2
%% type: 多选题
%% chapter: 磁场
%% point: 带电粒子在磁场中的运动
%% method: 无
%% score: 6
%% degree: 450
%% duplicateId: 0
%% body:
光滑刚性绝缘圆筒内存在着平行于轴的匀强磁场，筒上 P 点开有一个小孔，过 P 的横截面是以 O 为圆心的圆，如图所示。一带电粒子从 P 点沿 PO 射入，然后与筒壁发生碰撞。假设粒子在每次碰撞前、后瞬间，速度沿圆上碰撞点的切线方向的分量大小不变，沿法线方向的分量大小不变、方向相反；电荷量不变。不计重力。下列说法正确的是\xzanswer{BD}
\onepicture{figs/07.png}
\fourchoices[answer=BD]
{粒子的运动轨迹可能通过圆心 O}
{最少经 2 次碰撞，粒子就可能从小孔射出}
{射入小孔时粒子的速度越大，在圆内运动时间越短}
{每次碰撞后瞬间，粒子速度方向一定平行于碰撞点与圆心 O 的连线}
%% answer: BD
%% memo:
\memoanswer{
\onepicture[w=6cm]{figs/07_an.png}

D．假设粒子带负电，第一次从 A 点和筒壁发生碰撞如图，O$_{1}$ 为圆周运动的圆心。

由几何关系可知 $\angle$O$_{1}$AO 为直角，即粒子此时的速度方向为 OA，说明粒子在和筒壁碰撞时速度会反向，由圆的对称性在其它点撞击同理，D 正确；

A．假设粒子运动过程过 O 点，则过 P 点的速度的垂线和 OP 连线的中垂线是平行的不能交于一点确定圆心，由圆形对称性撞击筒壁以后的 A 点的速度垂线和 AO 连线的中垂线依旧平行不能确定圆心，则粒子不可能过 O 点，A 错误；

\onepicture[w=6cm]{figs/07_an2.png}

B．由题意可知粒子射出磁场以后的圆心组成的多边形应为以筒壁的内接圆的多边形，最少应为三角形如图所示，即撞击两次，B 正确；

C．速度越大粒子做圆周运动的半径越大，碰撞次数会可能增多，粒子运动时间不一定减少，C 错误。

故选 BD。
}

\item
%% number: 8
%% paperName: 2023年全国高考甲卷第 8 题 6 分
%% typeId: 2
%% type: 多选题
%% chapter: 电磁感应
%% point: 电磁感应中图象
%% method: 无
%% score: 6
%% degree: 450
%% duplicateId: 0
%% body:
一有机玻璃管竖直放在水平地面上，管上有漆包线绕成的线圈，线圈的两端与电流传感器相连，线圈在玻璃管上部的 5 匝均匀分布，下部的 3 匝也均匀分布，下部相邻两匝间的距离大于上部相邻两匝间的距离。如图（a）所示。现让一个很小的强磁体在玻璃管内沿轴线从上端口由静止下落，电流传感器测得线圈中电流 $I$ 随时间 $t$ 的变化如图（b）所示。则\xzanswer{AD}
\onepicture{figs/08.png}
\fourchoices[answer=AD]
{小磁体在玻璃管内下降速度越来越快}
{下落过程中，小磁体的 N 极、S 极上下顺倒了 8 次}
{下落过程中，小磁体受到的电磁阻力始终保持不变}
{与上部相比，小磁体通过线圈下部的过程中，磁通量变化率的最大值更大}
%% answer: AD
%% memo:
\memoanswer{
AD．电流的峰值越来越大，即小磁铁在依次穿过每个线圈的过程中磁通量的变化率越来越快，因此小磁体的速度越来越大，AD 正确；

B．假设小磁体是 N 极向下穿过线圈，则在穿入靠近每匝线圈的过程中磁通量向下增加，根据楞次定律可知线圈中产生逆时针的电流，而在穿出远离每匝线圈的过程中磁通量向下减少产生顺时针的电流，即电流方向相反，与题干图中描述的穿过线圈的过程电流方向变化相符，S 极向下同理；所以磁铁穿过 8 匝线圈过程中会出现 8 个这样的图像，并且随下落速度的增加，感应电流的最大值逐渐变大，所以磁体下落过程中磁极的 N、S 极没有颠倒，选项 B 错误；

C．线圈可等效为条形磁铁，线圈的电流越大则磁性越强，因此电流的大小是变化的小磁体受到的电磁阻力是变化的，不是一直不变的，C 错误。

故选 AD。
}

\item
%% number: 9
%% paperName: 2023年全国高考甲卷第 9 题 5 分
%% typeId: 4
%% type: 实验题
%% chapter: 电路
%% point: 伏安法测电阻
%% method: 无
%% score: 5
%% degree: 650
%% duplicateId: 0
%% body:
某同学用伏安法测绘一额定电压为 6 V、额定功率为 3 W 的小灯泡的伏安特性曲线，实验所用电压表内阻约为 6 k$\Omega$ 电流表内阻约为 1.5 $\Omega$。实验中有图（a）和（b）两个电路图供选择。
\begin{enumerate}
	\item
	实验中得到的电流 $I$ 和电压 $U$ 的关系曲线如图（c）所示，该同学选择的电路图是图\_\_\_\_\_\_\_\_（填“a”或“b”）
	\item
	若选择另一个电路图进行实验，在答题卡所给图上用实线画出实验中应得到的关系曲线的示意图\_\_\_\_\_\_。
\end{enumerate}
\onepicture[align=right]{figs/09.png}
%% answer: 
\jdanswer{
\begin{enumerate}
	\item
	a

	\item
	

	\item
	\onepicture[w=6cm]{figs/09_an.png}

\end{enumerate}
}
%% memo:
\memoanswer{
（1）由于 $\sqrt {{R_V}{R_A}}  \approx$ 95 $\Omega$，而小灯泡的电阻约为 $R = \frac{{{U^2}}}{P} \approx$ 12 $\Omega$，则电流表应采用外接法，则该同学选择的电路图是图（a）。

（2）若选用另一个电路图即图（b）实验，会有 $U = U_{\text{灯}} + IR_{A}$。则分别代入电流 200 mA、400 mA、500 mA，可知对应的电压应为 1.4 V、4.0 V、6.75 V，描点连线有
}

\item
%% number: 10
%% paperName: 2023年全国高考甲卷第 10 题 10 分
%% typeId: 4
%% type: 实验题
%% chapter: 直线运动
%% point: 匀变速直线运动
%% method: 无
%% score: 10
%% degree: 450
%% duplicateId: 0
%% body:
某同学利用如图（a）所示的实验装置探究物体做直线运动时平均速度与时间的关系。 让小车左端和纸带相连。右端用细绳跨过定滑轮和钩码相连。钩码下落，带动小车运动，打点计时器打出纸带。某次实验得到的纸带和相关数据如图（b）所示。

\begin{center}
\begin{tabular}{|c|}
\hline
位移区间 \\ \hline
\end{tabular}
\end{center}

AB

AC

AD

AE

AF

$\Delta x$（cm）

$\Delta x_{AD}$

$\overline v$（$\Ucms$）

$\overline v_{AD}$
\begin{enumerate}
	\item
	已知打出图（b）中相邻两个计数点的时间间隔均为 0.1 s．以打出 A 点时小车位置为初始位置，将打出 B、C、D、E、F 各点时小车的位移 $\Delta x$ 填到表中，小车发生该位移所用时间和平均速度分别为 $\Delta t$ 和 $\overline v$，表中 $\Delta x_{AD}$ = \_\_\_\_\_\_\_\_cm，$\overline v_{AD}$ = \_\_\_\_\_\_\_\_$\Ucms$。
	\item
	60
	\item
	60
	\item
	90
	\item
	30
	\item
	0
	\item
	0
	\item
	3
	\item
	6
	\item
	根据表中数据得到小车平均速度 $\overline v$ 随时间 $\Delta t$ 的变化关系，如图（c）所示。题卡上的图中补全实验点\_\_\_\_\_。
	\item
	从实验结果可知，小车运动的 $\overline v - \Delta t$ 图线可视为一条直线，此直线用方程 $\overline v = k\Delta t + b$ 表示，其中 $k$ = \_\_\_\_\_\_$\Ucms$$^{2}$，$b$ = \_\_\_\_\_\_\_\_\_$\Ucms$。（结果均保留 3 位有效数字）
	\item
	根据（3）中的直线方程可以判定小车做匀加速直线运动，得到打出 A 点时小车速度大小 $v_{A}$ = \_\_\_\_\_\_\_\_，小车的加速度大小 $a$ = \_\_\_\_\_\_\_\_。（结果用字母 $k$、$b$ 表示）
\end{enumerate}
\threepicture[labelA=2023全国甲卷10a, labelB=2023全国甲卷10b, labelC=2023全国甲卷10c, align=right]
{figs/10a.png}
{figs/10b.png}
{figs/10c.png}
%% answer: 
\jdanswer{
\begin{enumerate}
	\item
	24.00，80

	\item
	

	\item
	\onepicture[w=6cm]{figs/10_an.png}

	\item
	70.0，59.0

	\item
	$b$，2$k$

\end{enumerate}
}
%% memo:
\memoanswer{
（1）根据纸带的数据可得

$\Delta x_{AD} = \Delta x_{AB} +\Delta x_{BC} +\Delta x_{CD}$ = （6.60 + 8.00 + 9.40）cm = 24.00 cm

平均速度为 $\overline v_{AD} = \frac{{{x_{AD}}}}{{3T}}$ = 80 $\Ucms$

（2）根据第（1）小题结果补充表格和补全实验点图像得

\onepicture[w=6cm]{figs/10_an2.png}

（3）从实验结果可知，小车运动的 $\overline v - \Delta t$ 图线可视为一条直线，图像为

\onepicture[w=6cm]{figs/10_an3.png}

此直线用方程 $\overline v = k\Delta t + b$ 表示，由图像可知其中

$k = \frac{{101.0 - 59.0}}{{0.6}}$ $\Ucms$$^{2}$ = 70.0 $\Ucms$$^{2}$，$b$ = 59.0 $\Ucms$

（4）小球做匀变速直线运动，由位移公式 $x = v_{0}t + \frac{1}{2}at^{2}$，整理得 $\frac{x}{t} = v_{0} + \frac{1}{2}at$

即 $\overline v = v_{A} + \frac{1}{2}at$

故根据图像斜率和截距可得 $v_{A} = b$，$a$ = 2$k$
}

\item
%% number: 11
%% paperName: 2023年全国高考甲卷第 11 题 12 分
%% typeId: 6
%% type: 计算题
%% chapter: 曲线运动
%% point: 平抛运动
%% method: 无
%% score: 12
%% degree: 650
%% duplicateId: 0
%% body:
如图，光滑水平桌面上有一轻质弹簧，其一端固定在墙上。用质量为 $m$ 的小球压弹簧的另一端，使弹簧的弹性势能为 $E_{p}$。释放后，小球在弹簧作用下从静止开始在桌面上运动，与弹黄分离后，从桌面水平飞出。小球与水平地面碰撞后瞬间，其平行于地面的速度分量与碰撞前瞬间相等；垂直于地面的速度分量大小变为碰撞前瞬间的 $\frac{4}{5}$。小球与地而碰撞后，弹起的最大高度为 $h$。重力加速度大小为 $g$，忽略空气阻力。求
\begin{enumerate}
	\item
	小球离开桌面时的速度大小；
	\item
	小球第一次落地点距桌面上其飞出点的水平距离。
\end{enumerate}
\onepicture[align=right]{figs/11.png}
%% answer: 
\jdanswer{
\begin{enumerate}
	\item
	$\sqrt {\frac{{2{E_p}}}{m}} $

	\item
	$\frac{{5\sqrt {mgh{E_p}} }}{{2mg}}$

\end{enumerate}
}
%% memo:
\memoanswer{
（1）由小球和弹簧组成的系统机械能守恒可知

 $E_{p}=\frac{1}{2}mv^{2}$

得小球离开桌面时速度大小为  $v=\sqrt{\frac{2E_{p}}{m}}$

（2）离开桌面后由平抛运动规律可得

 $h=\frac{v_{y}^{\prime}}{2g}$

第一次碰撞前速度的竖直分量为  $v_{y}$，由题可知

 $v_{y}^{\prime}=\frac{4}{5}v_{y}$

离开桌面后由平抛运动规律得

 $x = vt, \quad V_{y} = gt$

解得小球第一次落地点距桌面上其飞出的水平距离为  $x = \frac{5}{2mg} mghE_{p}$
}

\item
%% number: 12
%% paperName: 2023年全国高考甲卷第 12 题 20 分
%% typeId: 6
%% type: 计算题
%% chapter: 运动和力
%% point: 动量和能量
%% method: 无
%% score: 20
%% degree: 450
%% duplicateId: 0
%% body:
如图，水平桌面上固定一光滑 U 型金属导轨，其平行部分的间距为 $l$，导轨的最右端与桌于右边缘对齐，导轨的电阻忽略不计。导轨所在区域有方向竖直向上的匀强磁场，磁感应强度大小为 $B$。一质量为 $m$、电阻为 $R$、长度也为 $l$ 的金属棒 P 静止在导轨上。导轨上质量为 3$m$ 的绝缘棒 Q 位于 P 的左侧，以大小为 $v_{0}$ 的速度向 P 运动并与 P 发生弹性碰撞，碰撞时间很短。碰撞一次后，P 和 Q 先后从导轨的最右端滑出导轨，并落在地面上同一地点。P 在导轨上运动时，两端与导轨接触良好，P 与 Q 始终平行。不计空气阻力。求
\begin{enumerate}
	\item
	金属棒 P 滑出导轨时的速度大小；
	\item
	金属体 P 在导轨上运动过程中产生的热量；
	\item
	与 P 碰撞后，绝缘棒 Q 在导轨上运动的时间。
\end{enumerate}
\onepicture[align=right]{figs/12.png}
%% answer: 
\jdanswer{
\begin{enumerate}
	\item
	$\frac{{2mR}}{{{B^2}{l^2}}}$

	\item
	$mv_{0}^{2}$

	\item
	$\frac{{{v_0}}}{2}$

\end{enumerate}
}
%% memo:
\memoanswer{
（1）由于绝缘棒 Q 与金属棒 P 发生弹性碰撞，根据动量守恒和机械能守恒可得  $3mv_{0}=3mv_{Q}+mv_{P}$

$$ \frac{1}{2}\times3mv_{0}^{2}=\frac{1}{2}\times3mv_{Q}^{2}+\frac{1}{2}mv_{P}^{2} $$

联立解得

$$ v_{\mathrm{P}}=\begin{array}{c}3v_{0}\\ 2\end{array},\quad v_{\mathrm{Q}}=\begin{array}{c}v_{0}\\ 2\end{array} $$

由题知，碰撞一次后，P 和 Q 先后从导轨的最右端滑出导轨，并落在地面上同一地点，则金属棒 P 滑出导轨时的速度大小为

 $v_{P}^{\prime}=v_{Q}=\frac{v_{0}}{2}$

（2）根据能量守恒有

 $\frac{1}{2}mv_{P}^{2}=\frac{1}{2}mv_{P}^{\prime2}+Q$

解得  $Q = mv_{0}^{2}$

（3）P、Q 碰撞后，对金属棒 P 分析，根据动量定理得

 $-B\bar{I}l\Delta t = mv_{\mathrm{P}}^{\prime} - mv_{\mathrm{P}}$

又  $q = \overline{I}\Delta t$， $\overline{I} = \frac{\overline{E}}{R} = \frac{\Delta\Phi}{R\Delta t} = \frac{Blx}{R\Delta t}$

联立可得  $x = \frac{mv_{0}R}{B^{2}l^{2}}$

由于 Q 为绝缘棒，无电流通过，做匀速直线运动，故 Q 运动的时间为

 $t = \frac{x}{v_{Q}} = \frac{2mR}{B^{2}l^{2}}$

}

\item
%% number: 13
%% paperName: 2023年全国高考甲卷第 13 题 5 分
%% typeId: 2
%% type: 多选题
%% chapter: 气体性质
%% point: 热力学第一定律
%% method: 无
%% score: 5
%% degree: 650
%% duplicateId: 0
%% body:
在一汽缸中用活塞封闭着一定量的理想气体，发生下列缓慢变化过程，气体一定与外界有热量交换的过程是\xzanswer{ABD}
\fivechoices[answer=ABD]
{气体的体积不变，温度升高}
{气体的体积减小，温度降低}
{气体的体积减小，温度升高}
{气体的体积增大，温度不变}
{气体的体积增大，温度降低}
%% answer: ABD
%% memo:
\memoanswer{
A．气体的体积不变温度升高，则气体的内能升高，体积不变气体做功为零，因此气体吸收热量，A 正确；

B．气体的体积减小温度降低，则气体的内能降低，体积减小外界对气体做功，由热力学第一定律 $\Delta U = W + Q$ 可知气体对外放热，B 正确；

C．气体的体积减小温度升高，则气体的内能升高，体积减小外界对气体做功，由热力学第一定律可知 $Q$ 可能等于零，即没有热量交换过程，C 错误；

D．气体的体积增大温度不变则气体的内能不变，体积增大气体对外界做功，由热力学第一定律可知 $Q$ > 0。即气体吸收热量，D 正确；

E．气体的体积增大温度降低则气体的内能降低，体积增大气体对外界做功，由热力学第一定律可知 $Q$ 可能等于零，即没有热量交换过程，E 错误。

故选 ABD。
}

\item
%% number: 13
%% paperName: 2023年全国高考甲卷第 13 题 10 分
%% typeId: 6
%% type: 计算题
%% chapter: 气体性质
%% point: 气体多过程
%% method: 无
%% score: 10
%% degree: 650
%% duplicateId: 0
%% body:
一高压舱内气体的压强为 1.2 个大气压，温度为 17 ${^\circ}C$，密度为 1.46 $\Ukgm$$^{3}$。
\begin{enumerate}
	\item
	升高气体温度并释放出舱内部分气体以保持压强不变，求气体温度升至 27 ${^\circ}C$ 时内气体的密度；
	\item
	保持温度 27 ${^\circ}C$ 不变，再释放出舱内部分气体使舱内压强降至 1.0 个大气压，求舱内气体的密度。
\end{enumerate}
%% answer: 
\jdanswer{
\begin{enumerate}
	\item
	1.41 $\Ukgm$$^{3}$

	\item
	1.18 $\Ukgm$$^{3}$

\end{enumerate}
}
%% memo:
\memoanswer{
A．气体的体积不变温度升高，则气体的内能升高，体积不变气体做功为零，因此气体吸收热量，A 正确；

B．气体的体积减小温度降低，则气体的内能降低，体积减小外界对气体做功，由热力学第一定律  $\Delta U = W + Q$ 可知气体对外放热，B 正确；

C．气体的体积减小温度升高，则气体的内能升高，体积减小外界对气体做功，由热力学第一定律可知 Q 可能等于零，即没有热量交换过程，C 错误；

D．气体的体积增大温度不变则气体的内能不变，体积增大气体对外界做功，由热力学第一定律可知 Q > 0。即气体吸收热量，D 正确；

E．气体的体积增大温度降低则气体的内能降低，体积增大气体对外界做功，由热力学第一

定律可知 Q 可能等于零，即没有热量交换过程，E 错误。

故选 ABD。

（10 分）一高压舱内气体的压强为 1.2 个大气压，温度为 17℃，密度为 1.46 $\Ukgm$ $^{3}$。

（1）升高气体温度并释放出舱内部分气体以保持压强不变，求气体温度升至  $27^{\circ}$C 时内气体的密度；

（2）保持温度  $27^{\circ}$C 不变，再释放出舱内部分气体使舱内压强降至 1.0 个大气压，求舱内气体的密度。

【详解】（1）由摄氏度和开尔文温度的关系可得

 $T_{1}=273+17\ K=290\ K,\ T_{2}=273+27\ K=300\ K$

理想气体状态方程  $pV = nRT$ 可知

 $nR=\frac{pV}{T}$

其中 n 为封闭气体的物质的量，即理想气体的  $\frac{pV}{T}$ 正比于气体的质量，则

 $\frac{\rho_{1}}{\rho_{2}}=\frac{\frac{m_{1}}{V}}{\frac{m_{2}}{V}}=\frac{\frac{p_{1}V}{T_{1}}}{\frac{p_{2}V}{T_{2}}}$

其中  $p_{1}=p_{2}=1.2p_{0},\rho_{1}=1.46\ \Ukgm^{3}$，代入数据解得

 $\rho_{2}=1.41\ \Ukgm^{3}$

（2）由题意得  $p_{3}=p_{0},\quad T_{3}=273+27\ K=300\ K$ 同理可得

 $\frac{\rho_{2}}{\rho_{3}}=\frac{\frac{m_{2}}{V}}{\frac{m_{3}}{V}}=\frac{\frac{p_{1}V}{T_{2}}}{\frac{p_{3}V}{T_{3}}}$

解得  $\rho_{3}=1.18\ \Ukgm^{3}$
}

\item
%% number: 14
%% paperName: 2023年全国高考甲卷第 14 题 5 分
%% typeId: 2
%% type: 多选题
%% chapter: 光学
%% point: 光的折射
%% method: 无
%% score: 5
%% degree: 650
%% duplicateId: 0
%% body:
等腰三角形 $\triangle$abc 为一棱镜的横截面，ab = ac；一平行于 bc 边的细光束从 ab 边射入棱镜，在 bc 边反射后从 ac 边射出，出射光分成了不同颜色的两束，甲光的出射点在乙光的下方，如图所示。不考虑多次反射。下列说法正确的是\xzanswer{ACE}
\onepicture{figs/14.png}
\fivechoices[answer=ACE]
{甲光的波长比乙光的长}
{甲光的频率比乙光的高}
{在棱镜中的传播速度，甲光比乙光的大}
{该棱镜对甲光的折射率大于对乙光的折射率}
{在棱镜内 be 边反射时的入射角，甲光比乙光的大}
%% answer: ACE
%% memo:
\memoanswer{
ABD．根据折射定律和反射定律作出光路图如图所示

\onepicture[w=6cm]{figs/14_an.png}

由图可知，乙光的折射角较小，根据折射定律可知乙光的折射率大，则乙光的频率大，根据 $c = f\lambda$ 可知，乙光的波长短，A 正确、BD 错误；

C．根据 $v = \frac{c}{n}$ 可知在棱镜中的传播速度，甲光比乙光的大，C 正确；

E．根据几何关系可知光在棱镜内 bc 边反射时的入射角，甲光比乙光的大，E 正确。

故选 ACE
}

\item
%% number: 14
%% paperName: 2023年全国高考甲卷第 14 题 10 分
%% typeId: 6
%% type: 计算题
%% chapter: 机械波
%% point: 干涉衍射
%% method: 无
%% score: 10
%% degree: 550
%% duplicateId: 0
%% body:
分别沿 $x$ 轴正向和负向传播的两列简谐横波 P、Q 的振动方向相同振幅均为 5 cm，波长均为 8 m，波速均为 4 $\Ums$。$t$ = 0 时刻，P 波刚好传播到坐标原点，该处的质点将自平衡位置向下振动；Q 波刚好传到 $x$ = 10 m 处，该处的质点将自平衡位置向上振动。经过一段时间后，两列波相遇。
\begin{enumerate}
	\item
	在答题卡给出的坐标图上分别画出 P、Q 两列波在 $t$ = 2.5 s 时刻的波形图（P 波用虚线，Q 波用实线）；
	\item
	求出图示范围内的介质中，因两列波干涉而振动振幅最大和振幅最小的平衡位置。
\end{enumerate}
\onepicture[align=right]{figs/14.png}
%% answer: 
\jdanswer{
\begin{enumerate}
	\item
	

	\item
	\onepicture[w=6cm]{figs/14_an2.png}

	\item
	振幅最大的平衡位置有 $x$ = 3 m、$x$ = 7 m；

	\item
	振幅最小的平衡位置有 $x$ = 1 m、$x$ = 5 m、$x$ = 9 m。

\end{enumerate}
}
%% memo:
\memoanswer{
ABD. 根据折射定律和反射定律作出光路图如图所示

由图可知，乙光的折射角较小，根据折射定律可知乙光的折射率大，则乙光的频率大，根据  $c = f\lambda$ 可知，乙光的波长短，A 正确、BD 错误；

C. 根据  $v = \frac{c}{n}$ 可知在棱镜中的传播速度，甲光比乙光的大，C 正确；

E. 根据几何关系可知光在棱镜内 bc 边反射时的入射角，甲光比乙光的大，E 正确。故选 ACE。

（10 分）分别沿 x 轴正向和负向传播的两列简谐横波 P、Q 的振动方向相同振幅均为 5 cm，波长均为 8 m，波速均为 4 $\Ums$。t = 0 时刻，P 波刚好传播到坐标原点，该处的质点将自平衡位置向下振动；Q 波刚好传到 x = 10 m 处，该处的质点将自平衡位置向上振动。经过一段时间后，两列波相遇。

（1）在答题卡给出的坐标图上分别画出 P、Q 两列波在  $t=2.5\ s$ 时刻的波形图（P 波用虚线，Q 波用实线）；

（2）求出图示范围内的介质中，因两列波干涉而振动振幅最大和振幅最小的平衡位置。

【详解】（1）根据  $\Delta x = vt$ 得

$\Delta x = 4 \times 2.5 \, m = 10 \, m$

可知  $t = 2.5 \, s$ 时 P 波刚好传播到  $x = 10 \, m$ 处，Q 波刚好传播到  $x = 0$ 处，根据上坡下坡法可得波形图如图所示

（2）根据题意可知，P、Q 两波振动频率相同，振动方向相反，两波叠加时，振动加强点

的条件为到两波源的距离差  $\Delta x = \frac{(2n+1)\lambda}{2}(n=0,1,2)$

解得振幅最大的平衡位置有  $x = 3 \, m$、 $x = 7 \, m$

振动减弱的条件为  $\Delta x = n\lambda (n = 0, 1, 2)$

解得振幅最小的平衡位置有  $x = 1 \, m$、 $x = 5 \, m$、 $x = 9 \, m$
}

\end{enumerate}

\end{document}
